\section{Discussion}
\label{sec:discussion}

The experiments form a hierarchy of evidence. Spectral visibility says where
to look, behavior-labeled capture asks whether a candidate overlaps a known
axis, matched controls identify which training condition the direction tracks,
and interventions test whether the model is sensitive to those coordinates.
No single stage supports the full claim on its own.

\paragraph{The spectrum is a search device, not an explanation.}
Behavior-labeled activations can identify a refusal direction directly
\citep{arditi2024}. Our spectral analysis asks whether the Base-to-Instruct
checkpoint delta concentrates change in the same subspace. Prompt-labeled capture
and the spectral-versus-random intervention gap show that measured refusal is
sensitive to the selected leading subspace. The additive model supplies an ideal outlier
threshold \citep{benaych2012}; its fitted empirical counterpart
prioritizes directions for those tests. The practical role of the spectrum is
therefore compression: it turns a high-dimensional endpoint delta into a small
set of hypotheses that can be challenged with behavioral evidence.

\paragraph{What the controls establish.}
The controls separate four claims:
\begin{itemize}
\setlength{\itemsep}{1pt}
\item \textbf{MP and Gaussian references} measure departure from an
iid-isotropic benchmark, not alignment specificity.
\item \textbf{A seeded random subspace} controls for projection dimension in the
top-$128$ audit, but does not estimate a null distribution.
\item \textbf{Matched benign controls} isolate the joint target-policy and content
contrast within the medical organism.
\item \textbf{Interventions} show sensitivity in the tested models, not circuit
locality, one-dimensional sufficiency, or prospective prediction.
\end{itemize}
Together, these controls support a workflow rather than a universal safety
vector: use geometry to nominate directions, matched contrasts to attach
condition information, and intervention to test the resulting hypothesis.
The unrestricted delta may reflect generic fine-tuning, optimizer correlations,
or low-dimensional adaptation \citep{aghajanyan2021,hu2022lora}; matched arm
labels supply the condition information. Across three medical organisms the
direction recurs, but raw update magnitude remains unmatched.

\paragraph{Limitations.}
The refusal census covers one released 8B model \citep{grattafiori2024llama3},
and the 14B organism reproduces geometry but not causal advantage over random
ablation. Each endpoint delta aggregates training stages, while each organism
fixes one recipe.

Global projection does not localize a circuit, and each medical projection uses
a checkpoint that contributed to its fitted direction. Leave-one-pair-out is a
ranking screen, not a held-out causal intervention. The refusal direction may encode prompt
distribution, topic, or style alongside refusal. HarmBench reproduces the spectral-versus-random gap on
held-out harmful prompts \citep{mazeika2024harmbench}. Projection is blunt and
damages MMLU, ARC-C, and GSM8K
\citep{hendrycks2021mmlu,cobbe2021gsm8k,clark2018arc}; steering is nonmonotonic.

The organism claim rests on matched separation, held-out ranking, and
cross-family recurrence, not prevalence. Baselines do not favor weight-SVD:
row-mean and weight-SVD exchange margin order under score squaring, activation
PCA ties their arm ordering, folds overlap, and update magnitude remains
confounded. The code-organism and 14B causal audits fail their frozen criteria.
Prospective use requires directions and
thresholds frozen before new endpoint deltas are observed.

\section*{Ethical considerations}
The refusal intervention demonstrates a method for weakening a safety-related
behavioral score. It also causes severe capability damage and does not provide a
capability-preserving bypass, but the same analysis could still inform attempts
to remove safeguards. We therefore present the intervention as a diagnostic
result rather than a deployment recipe. The medical model-organism outputs are
synthetic and are not medical advice. Their local-judge labels have not been
independently or human validated. Any operational use would require stronger behavioral
evaluation, human review, held-out causal tests, and explicit controls for
capability and misuse risk.
